% Lab 1: Orientation, set-up, and your first app
% Application Development for Mobile Devices
\documentclass[10pt,aspectratio=169,t]{beamer}
\usetheme[lab]{upbminimal}

\course[ADMD]{Application Development for Mobile Devices}
\decklabel{Lab 1}
\title{Orientation, set-up, and your first app}
\subtitle{Environment, tooling and a first app on a device}
\author{Dinu-Ștefan Rusu}
\date{Week 2}

\begin{document}

\titleframe

\begin{frame}[eyebrow=Today]{Lab objectives}
  \vskip 2mm
  \begin{grid}[2]
    \cell[\faWrench]{Install a working toolchain}{Flutter and Dart on your laptop, with flutter doctor reporting no blocking issues.}
    \cell[\faServer]{Run an app on a device}{An emulator or simulator that boots, with the default Flutter app running on it.}
  \end{grid}
  \vskip 6mm
  \taskmeta{Time}{2 hours}{Present}{To the lab instructor, at the end of this lab}
\end{frame}

\begin{frame}[eyebrow=Setup]{The cumulative lab app}
  \begin{steps}
    \item One Flutter project per student. Every lab continues it.
    \item At the end of each lab, present what you built to the lab instructor.
    \item Install the tools before starting.
  \end{steps}
  \vskip 4mm
  \begin{takeaway}{Seven labs build one app that grows from empty to networked and authenticated.}
    Lab 1 sets up the machine. Lab 7 is the test. Labs 2 to 6 add features.
  \end{takeaway}
\end{frame}

\begin{frame}[eyebrow=Grading]{Grading}
  {\usebeamerfont{small}
  \begin{tabular}[t]{lp{6cm}r}
    \thead{Work} & \thead{Description} & \thead{Points} \\
    Colloquium & End of semester, before the exam session & 20 \\
    Midterm & Middle of the semester & 20 \\
    Lab assignments & Labs 1 to 6, cumulative app & 30 \\
    Practical test & Lab 7, 90 minutes, individual & 30 \\
    \midrule
    & \textbf{Pass mark} & \textbf{50 of 100} \\
  \end{tabular}
  }
  \vskip 2mm
  \begin{hint}
    Each lab is graded when you present it. Be ready to run the app and explain any line.
  \end{hint}
\end{frame}

\begin{frame}[fragile,eyebrow=Setup]{Installation links}
  \begin{columns}[T]
    \begin{column}{0.55\textwidth}
\begin{shell}
# Flutter SDK:
# https://docs.flutter.dev/install
# Add flutter/bin to your PATH

flutter --version
flutter doctor -v

# Android Studio or Xcode
# Pick one IDE: VS Code or IntelliJ
\end{shell}
    \end{column}
    \begin{column}{0.41\textwidth}
      \kv{Before starting}{}
      Install Flutter, an IDE, and either Android Studio or Xcode.
      \\[2mm]
      \kv{Target}{}
      The current stable Flutter channel.
      \\[2mm]
      \kv{Important}{}
      Type URLs exactly. They are complete on one line.
    \end{column}
  \end{columns}
\end{frame}

\begin{frame}[eyebrow=Task I]{Environment setup}
  \begin{columns}[T]
    \begin{column}{0.55\textwidth}
      \begin{steps}
        \item Install Flutter from \link{https://docs.flutter.dev/install}{docs.flutter.dev/install}.
        \item Follow Install manually for your OS.
        \item Add flutter/bin to your PATH.
        \item Install Android Studio or Xcode (macOS).
        \item Install an editor: VS Code or IntelliJ.
        \item Run \code{flutter doctor}, fix every blocking issue.
      \end{steps}
    \end{column}
    \begin{column}{0.41\textwidth}
      \kv{Done when}{}
      \begin{checklist}
        \item \code{flutter --version} works in a new terminal.
        \item \code{flutter doctor} shows no red warnings for Flutter or Android.
        \item You can explain one warning you fixed.
      \end{checklist}
    \end{column}
  \end{columns}
\end{frame}

\begin{frame}[eyebrow=Task II]{Running the default app}
  \begin{columns}[T]
    \begin{column}{0.55\textwidth}
      \begin{steps}
        \item Boot an emulator or simulator.
        \item Run \code{flutter devices} to confirm it appears.
        \item Create the project: \code{flutter create lab1\_app}.
        \item Run it: \code{flutter run}.
        \item Change the app title string in \code{lib/main.dart}.
        \item Press r to hot-reload and see it update.
      \end{steps}
    \end{column}
    \begin{column}{0.41\textwidth}
      \kv{Done when}{}
      \begin{checklist}
        \item The counter app runs and the button increments it.
        \item Your edited title appears after hot reload.
        \item You can repeat the hot reload on request.
      \end{checklist}
    \end{column}
  \end{columns}
\end{frame}

\begin{frame}[fragile,eyebrow=Troubleshooting]{flutter doctor errors}
  \trouble{cmdline-tools missing}{Android Studio > SDK Manager > SDK Tools > Android SDK Command-line Tools. Install, then run \code{flutter doctor} again.}
  \trouble{License unknown}{Run \code{flutter doctor --android-licenses} and answer \code{y} to all prompts.}
  \trouble{Unable to locate SDK}{Open Android Studio once. The SDK downloads on first launch. On Mac, accept the Xcode license first.}
\end{frame}

\begin{frame}[fragile,eyebrow=Troubleshooting]{Emulators and devices}
  \trouble{Emulator never boots}{Enable hardware acceleration (VT-x / AMD-V) in BIOS. On Apple silicon Macs, pick arm64 images only.}
  \trouble{No devices listed}{Run \code{flutter devices}. Start the emulator with \code{flutter emulators --launch <id>}. Physical Android phones work with USB debugging on.}
  \trouble{Xcode but no simulator}{Open Xcode once and accept the license, then run \code{open -a Simulator}. If it still fails, run \code{xcode-select --install}.}
\end{frame}

\begin{frame}[eyebrow=GitHub]{Keeping your work in Git}
  \begin{steps}
    \item Create a repository on GitHub.
    \item Clone it: \code{git clone <url>}.
    \item Work on your code, commit often with clear messages.
    \item Push your commits: \code{git push}.
  \end{steps}
  \vskip 3mm
  \muted{The app grows across seven labs. A commit after each task keeps every working state one command away.}
\end{frame}

\begin{frame}[eyebrow=Presentation]{Presenting your work}
  \begin{columns}[T]
    \begin{column}{0.55\textwidth}
      \begin{steps}
        \item Call the lab instructor when your checklists are done, or in the last twenty minutes.
        \item Run the app live, on your emulator or device.
        \item Walk through each task and answer questions on your code.
      \end{steps}
    \end{column}
    \begin{column}{0.41\textwidth}
      \kv{Today}{}
      Show \code{flutter doctor}, the counter app running, and your edited title after a hot reload.
    \end{column}
  \end{columns}
  \begin{takeaway}{A lab can only be presented in its own session.}
  \end{takeaway}
\end{frame}

\closingframe{Presentation}{The app running · one hot reload, shown live}

\end{document}
